\documentclass[10pt,a4paper]{amsart}
\usepackage[margin=19mm]{geometry}
\usepackage{amsmath,amssymb,mathtools}
\usepackage[scr=rsfs,cal=euler]{mathalfa}
\usepackage{xfrac}
\usepackage[unicode,hidelinks,bookmarks=false]{hyperref}
\newcommand{\ie}{i.e.\ }
\newcommand{\cf}{cf.\ }
\newcommand{\resp}{respectively\ }
\newcommand{\Subsection}{\subsection}
\providecommand{\nobreakdash}{\nobreak\hbox{-}}
\setlength{\emergencystretch}{2em}
\multlinegap0pt
\usepackage{xcolor}
\usepackage{float}
\usepackage{amssymb}
\usepackage[normalem]{ulem}
\usepackage{dsfont}
\usepackage{enumerate}
\usepackage{tikz}
\newtheorem{lemma}{Lemma}
\newtheorem{theorem}{Theorem}
\def\L{\mathcal{L}}
\newcommand\cT{{\mathcal T}}
\newcommand{\eps}{\varepsilon}
\title{Fixed-point approximation for\\ self-consistent transfer operators with Newton's method}
\author{Wael Bahsoun, Gary Froyland, Maxence Phalempin}
\date{Source: arXiv:2605.08803v1 (CC BY 4.0)}
\begin{document}
\maketitle

\noindent\emph{Source and licence.} Fixed-point approximation for\\ self-consistent transfer operators with Newton's method. Version arXiv:2605.08803v1 (CC BY 4.0), distributed by arXiv under the Creative Commons Attribution 4.0 International licence, http://creativecommons.org/licenses/by/4.0/, as recorded in the arXiv licence metadata for this exact version and shown on the abstract page https://arxiv.org/abs/2605.08803v1. Full licence notice: \texttt{licenses/arxiv-2605.08803-CC-BY-4.0.txt}.\par\smallskip

\noindent\textit{Editorial scope.} Counted unit: the article's lemma lem:convhN (source0.tex lines 370--373). The article's complete original proof of it is delivered with it (source0.tex lines 375--453, one \texttt{proof} environment of 79 lines). No mathematics is written, repaired or re-derived: the statement and the proof are the article's own lines, in the article's own notation, and no retained line is substituted or re-flowed.  Retained uncounted as the source-local context the proof needs: lines 136--138; lines 250--256; lines 308--323; lines 336--341; lines 1476--1489; lines 1481--1483; lines 1485--1487; lines 1496--1516.  Nothing was omitted from any retained span, so the package carries no omission record.  No retained line contains a citation, so the wrapper carries no bibliography and no bibliography entry.  No retained line contains a figure environment, an image-inclusion command or drawing code, so no image is delivered and the package declares no asset dependency.  Editorial formatting only, in addition: an AMS \texttt{amsart} preamble that restates the article's own declarations and macros used by the retained lines, namely the environments \texttt{lemma}, \texttt{theorem} on separate counters, together with the standard packages those lines need.  The retained environments print in this document as Lemma 1, Theorem 1 in order of appearance, whereas the article prints the same items with the numbering of its own full document.  The complete native source of the article as distributed in the arXiv e-print for this exact version is bundled unchanged as \texttt{sources/200176/source0.tex}.
\begin{equation}\label{eq:evolve}
h^{n+1}=\L_{T_{\eps,h^{n}}}h^n.
\end{equation}

\begin{lemma}\label{lemkern}
The following statements hold:
\begin{itemize}
\item[1)] for $f \in W^{i,1}$, $i=0,1,\dots$, $\|\Pi_Nf\|_{W^{i,1}}\leq \|f\|_{W^{i,1}}$;
\item[2)] $\exists C>0$ such that for $f \in W^{i+1,1}$, $i=0,1,\dots$, $$\|\Pi_N f-f\|_{W^{i,1}}\leq C\frac{\ln(N)}{N}\|f\|_{W^{i+1,1}}.$$
\end{itemize} 
\end{lemma}

\begin{lemma}\label{lem:discLY_co}
For any $N\in\mathbb N$ the operator $\Pi_N\tilde\L_\eps$ satisfies the following:
\begin{enumerate}
\item For any $\eps\in[0,\eps^*_0]$, for all $h\in W^{i,1}$, $i=1,2,3,4$
\begin{equation}
\|\Pi_N\tilde\L_\eps h\|_{W^{i,1}}\leq \lambda^{i}\|h\|_{W^{i,1}}+ C_{LY}\|h\|_{W^{i-1,1}},\nonumber
\end{equation}
where $C_{LY}$ and $\lambda$ are the same constants as in \eqref{eq:seqLY}.
Moreover, exists $\tilde C_{LY}>0$ such that for and any sequence  $(f_j)_{j\in \mathbb{N}} 
\in(\mathcal D_1)^{\mathbb N}$,
$$
 \|\Pi_N\L_{T_{\eps,f_n}}\circ\cdots \circ \Pi_N\L_{T_{\eps, f_1}}h\|_{W^{i,1}}\le \lambda^{in}\|h\|_{W^{i,1}}+\tilde C_{LY}\|h\|_{W^{i-1,1}}.
$$
\item for $\eps\in[0,\eps_0^*]$, $\Pi_N\tilde\L_\eps$ leaves $B_K$ invariant and hence has a fixed point $h_N^*\in B_K$. 
\end{enumerate}
\end{lemma}

\begin{lemma}\label{lem:contracto}
There are $\eps_1^*>0$ and $N_0\in \mathbb N$, $\tilde \theta (N_0) >0$, such that for any $N\geq N_0$, $n\in\mathbb{N}$, $\eps\in[0,\eps_1^*]$, $h\in V^{i,1}$, $i=1,2$ and any sequence  $(f_j)_{j\in \mathbb{N}} \in (B_K)^{\mathbb{N}}$,
\begin{align*}
    \|\Pi_N\L_{T_{\eps,f_n}}\circ\cdots \circ \Pi_N\L_{T_{\eps, f_1}}h\|_{W^{i,1}}\le C \tilde \theta^n\|h\|_{W^{i,1}}.
\end{align*}
\end{lemma}

\begin{lemma}
    \label{introprop}
The following sequential Lasota-Yorke and continuity inequalities hold\footnote{In the following two inequalities we label different elements of $\mathcal{D}_1$ by a subscript; i.e., $f_j$ for some $j$. This should not be confused with the superscript notation used below in \eqref{eq:evolve}, $h^n$, for a state at time $n$.}: 
\begin{enumerate}
\item $\exists C_{LY}>0$ such that for all $T_{\eps,f_n},\dots,T_{\eps,f_1}\in\cT$, $h\in W^{i,1}$, $i=1,2,3,4$, and $n\in\mathbb N$
\begin{equation}\label{eq:seqLY}
\|\L_{ T_{\eps,f_n}\circ\dots\circ T_{\eps,f_1}} h\|_{W^{i,1}}\leq \lambda^{ni}\|h\|_{W^{i,1}}+ C_{LY}\|h\|_{W^{i-1,1}};
\end{equation}
\item $\exists C>0$ such that for all $T_{\eps,f_2},T_{\eps,f_1}\in\cT$, $f_1,f_2\in W^{i,1}$, $h\in W^{i+1,1}$, $i=0,1$,
\begin{equation}\label{eq:closeop}
\|(\L_{T_{\eps, f_1}}-\L_{T_{\eps,f_2}})h\|_{W^{i,1}} \leq  C\eps \|f_1-f_2\|_{W^{i,1}} \|h\|_{W^{i+1,1}}.
\end{equation}
\end{enumerate}
\end{lemma}

\begin{equation}
\|\L_{ T_{\eps,f_n}\circ\dots\circ T_{\eps,f_1}} h\|_{W^{i,1}}\leq \lambda^{ni}\|h\|_{W^{i,1}}+ C_{LY}\|h\|_{W^{i-1,1}};
\end{equation}

\begin{equation}
\|(\L_{T_{\eps, f_1}}-\L_{T_{\eps,f_2}})h\|_{W^{i,1}} \leq  C\eps \|f_1-f_2\|_{W^{i,1}} \|h\|_{W^{i+1,1}}.
\end{equation}

\begin{theorem}\label{thm:gen}
 The following hold: 
\begin{enumerate}
    \item For each
    $\eps\in[0,\eps^*_0]$, $\tilde\L_\eps$ has a fixed point in $B_K$.    
\item There exists $\eps_1^* \in(0,\eps_0^*]$, $C>0$, $\theta\in(0,1)$ such that for all $\eps \in[0,\eps^*_1]$, $h\in W^{i,1}$, $i=1,2$, with $\int h=0$ and any sequence  $(f_j)_{j\in \mathbb{N}} \in (B_K)^{\mathbb{N}}$,
\begin{align*}
    \|\L_{T_{\eps,f_n}}\circ\cdots \circ \L_{T_{\eps, f_1}}h\|_{W^{i,1}}\le C \tilde \theta^n\|h\|_{W^{i,1}}.
\end{align*}




\item There exists $\eps_2^* \in(0,\eps_1^*]$,  such that for all $\eps \in[0,\eps^*_2] $ the operator $\tilde\L_\eps$ has a unique fixed point $h_\eps^*\in B_K$. Moreover, there exists $C_K>0$ and $\gamma\in (0,1)$ such that for all $\eps \in[0,\eps^*_2]$,  for all $h^0\in B_K$, we have
\[
\|{\tilde\L}_\eps^{n}(h^0)-h_\eps^*\|_{W^{1,1}}\le C_K \gamma^n;
\]
more explicitly,  
$$\|\mathcal{L}_{T_{\eps, h^{n-1}}}\circ\cdots\circ \mathcal{L}_{T_{\eps, h^{1}}}\circ\mathcal{L}_{T_{\eps, h^{0}}}h^{0} - h^*_\eps\|_{L^1}\le C_K \gamma^n.$$
\end{enumerate}
\end{theorem}

\begin{lemma}\label{lem:convhN}
There is $\eps_{3}^*\in (0, \eps_2^*)$ such that for $\eps\in [0,\eps_{3}^*]$, if $h_\varepsilon^*$ is the unique fixed point of $\tilde\L_\eps$ and $h_{\varepsilon,N}^*$ is a fixed point of  $\Pi_N\tilde\L_\eps$ then
$$\|h_{\varepsilon,N}^*-h_\eps^*\|_{W^{1,1}}=O\left(\frac{\ln N}{N}\right).$$
\end{lemma}

\begin{proof}

For $\eps\in[0,\eps^*_{3}]$, for any $N\in\mathbb N$, by Lemma \ref{lem:discLY_co}, the operator $\Pi_N\tilde\L_\eps$ has a fixed point $h_{\eps,N}^*\in B_K$, and by Theorem \ref{thm:gen},
\[
\|{\tilde\L}_\eps^{n}(h_{\eps,N}^*)-h_\eps^*\|_{W^{1,1}}\le C_K \gamma^n,
\]
where $h_\eps^*\in B_K$ is the unique fixed of $\tilde\L_\eps$. We have
\begin{equation}\label{eq:fixpoint}
\begin{split}
\|h_{\varepsilon,N}^*-h_\eps^*\|_{W^{1,1}}\le& \|(\Pi_N\tilde\L_\eps)^n(h_{\varepsilon,N}^*)-{\tilde\L}_\eps^{n}(h_{\eps,N}^*)\|_{W^{1,1}}+\|{\tilde\L}_\eps^{n}(h_{\eps,N}^*)-h_\eps^*\|_{W^{1,1}}\\
&\le\|(\Pi_N\tilde\L_\eps)^n(h_{\varepsilon,N}^*)-{\tilde\L}_\eps^{n}(h_{\eps,N}^*)\|_{W^{1,1}} +C_K \gamma^n.
\end{split}
\end{equation}
Where $(\Pi_N\tilde\L_\eps)^n$ is a notation for the $n^{th}$ iterate of the map $\Pi_N\tilde\L_\eps$.  Define the sequence $(h_{\eps,N}^{*_{n}})_{n\in \mathbb N}$ by $h_{\eps,N}^{*_n}:=\tilde \L_\eps^n h_N^*$. Then,
\begin{align*}
&\|(\Pi_N\tilde\L_\eps)^n(h_{\varepsilon,N}^*)-{\tilde\L}_\eps^{n}(h_{\eps,N}^*)\|_{W^{1,1}}\\
&=\|\sum_{k=0}^{n-1}(\Pi_N\circ \L_{T_{\eps,h_{\varepsilon,N}^*}})^{n-k}\left(\Pi_N \L_{\eps,T_{\eps,h_N^*}}-\L_{T_{\eps,h_{\eps,N}^{*_k}}}\right) \L_{T_{\eps,h_{\eps,N}^{*_{k-1}}}}\circ\cdots \circ\L_{T_{\eps, h_{\eps,N}^{*}}}h_{\eps,N}^*\|_{W^{1, 1}}\nonumber\\
&=\|\sum_{k=0}^{n-1}(\Pi_N\circ \L_{T_{\eps,h_{\varepsilon,N}^*}})^{n-k}\left[\left(\Pi_N-I\right) \L_{\eps,T_{\eps,h_N^*}}+\left(\L_{\eps,T_{\eps,h_N^*}}-\L_{T_{\eps,h_{\eps,N}^{*_k}}}\right)\right] \L_{T_{\eps,h_{\eps,N}^{*_{k-1}}}}\circ\cdots \circ\L_{T_{\eps, h_{\eps,N}^{*}}}h_{\eps,N}^*\|_{W^{1, 1}}\nonumber\\
&=(I)+(II),
    \end{align*}
where 
$$(I):= \sum_{k=0}^{n-1}(\Pi_N\circ \L_{T_{\eps,h_{\varepsilon,N}^*}})^{n-k}\left[\left(\Pi_N-I\right) \L_{\eps,T_{\eps,h_N^*}}\right] \L_{T_{\eps,h_{\eps,N}^{*_{k-1}}}}\circ\cdots \circ\L_{T_{\eps, h_{\eps,N}^{*}}}h_{\eps,N}^*\|_{W^{1, 1}},$$
and 
$$
(II):=\|\sum_{k=0}^{n-1}(\Pi_N\circ \L_{T_{\eps,h_{\varepsilon,N}^*}})^{n-k}\left[\L_{\eps,T_{\eps,h_N^*}}-\L_{T_{\eps,h_{\eps,N}^{*_k}}}\right] \L_{T_{\eps,h_{\eps,N}^{*_{k-1}}}}\circ\cdots \circ\L_{T_{\eps, h_{\eps,N}^{*}}}h_{\eps,N}^*\|_{W^{1, 1}}.
$$
We first treat $(I)$. Note in $(I)$, $(\Pi_N\circ \L_{T_{\eps,h_{\varepsilon,N}^*}})^{n-k}$ acts on a zero-mean function, thus Lemma \ref{lem:contracto} applies and using Lemma \ref{lemkern} alongside it to control the term $(\Pi_N-I)$, we obtain:

\begin{align}
(I)&=\|\sum_{k=0}^{n-1}(\Pi_N\circ \L_{T_{\eps,h_{\varepsilon,N}^*}})^{n-k}\left[\left(\Pi_N-I\right) \L_{\eps,T_{\eps,h_N^*}}\right] \L_{T_{\eps,h_{\eps,N}^{*_{k-1}}}}\circ\cdots \circ\L_{T_{\eps, h_{\eps,N}^{*}}}h_{\eps,N}^*\|_{W^{1, 1}}\nonumber\\
&\leq \sum_{k=0}^{n-1}\tilde \theta^{n-k}\|\left(\Pi_N-I\right)\|_{W^{2,1}\to W^{1,1}} \|\L_{\eps,T_{\eps,h_N^*}}\circ \L_{T_{\eps,h_{\eps,N}^{*_{k-1}}}}\circ\cdots \circ\L_{T_{\eps, h_{\eps,N}^{*}}}h_{\eps,N}^*\|_{W^{2, 1}}\nonumber\\
&\leq C\frac{\ln(N)}{N}\sum_{k=0}^{n-1}\tilde \theta^{n-k} \|\L_{\eps,T_{\eps,h_N^*}}\circ \L_{T_{\eps,h_{\eps,N}^{*_{k-1}}}}\circ\cdots \circ\L_{T_{\eps, h_{\eps,N}^{*}}}h_{\eps,N}^*\|_{W^{2, 1}}\nonumber\\
&\leq KC\frac{\ln(N)}{N}\sum_{k=0}^{n-1}\tilde \theta^{n-k},\label{eq:fixpartun}
\end{align}
where in the last line we used that $h_N^*\in B_K$ and Lemma \ref{lem:discLY_co}.

We now deal with $(II)$, using again Lemma \ref{lem:contracto} and \eqref{eq:closeop} from Lemma \ref{introprop},
\begin{align}
(II)&=\|\sum_{k=0}^{n-1}(\Pi_N\circ \L_{T_{\eps,h_{\varepsilon,N}^*}})^{n-k}\left[\L_{\eps,T_{\eps,h_N^*}}-\L_{T_{\eps,h_{\eps,N}^{*_k}}}\right] \L_{T_{\eps,h_{\eps,N}^{*_{k-1}}}}\circ\cdots \circ\L_{T_{\eps, h_{\eps,N}^{*}}}h_{\eps,N}^*\|_{W^{1, 1}}\nonumber\\
&\leq \sum_{k=0}^{n-1}\tilde \theta^{n-k}\|\left[\L_{\eps,T_{\eps,h_N^*}}-\L_{T_{\eps,h_{\eps,N}^{*_k}}}\right] \L_{T_{\eps,h_{\eps,N}^{*_{k-1}}}}\circ\cdots \circ\L_{T_{\eps, h_{\eps,N}^{*}}}h_{\eps,N}^*\|_{W^{1, 1}}\nonumber\\
&\leq \sum_{k=0}^{n-1}\tilde \theta^{n-k} \eps \|h_N^*-h_{\eps,N}^{*_k}\|_{W^{1,1}} \|\L_{T_{\eps,h_{\eps,N}^{*_{k-1}}}}\circ\cdots \circ\L_{T_{\eps, h_{\eps,N}^{*}}}h_{\eps,N}^*\|_{W^{2, 1}}\nonumber\\
&\leq K\eps\sum_{k=1}^{n-1}\tilde \theta^{n-k}\|h_N^*-h_{\eps,N}^{*_k}\|_{W^{1,1}},\label{eq:fixpartdeux}
\end{align}
where the last line is obtained using the fact that $h_N^*\in B_K$ and \eqref{eq:seqLY}.
Let $$U_n^N:=\|(\Pi_N\tilde\L_\eps)^n(h_{\varepsilon,N}^*)-{\tilde\L}_\eps^{n}(h_{\eps,N}^*)\|_{W^{1,1}}.$$ 
Notice that by \eqref{eq:fixpartun} and \eqref{eq:fixpartdeux}, satisfied by $(I)$ and $(II)$, we obtain the following inequality

\begin{align*}
U_n^N &\leq KC\frac{\ln(N)}{N}\sum_{k=0}^{n-1}\tilde \theta^{n-k}+K\eps\sum_{k=0}^{n-1}\tilde \theta^{n-k}U_k^N\\
&\leq KC\frac{\ln(N)}{N}\frac{\tilde \theta}{1-\tilde \theta}+K\eps\sum_{k=0}^{n-1}\tilde \theta^{n-k}U_k^N.
\end{align*}
Therefore, for any $n\in \mathbb N$, $U_n^N\leq V_n^N$ where $(V_n^N)_{n\in \mathbb N}$ is a sequence defined inductively by

\begin{align*}
  V_n^N=KC\frac{\ln(N)}{N}\frac{\tilde \theta}{1-\tilde \theta}+K\eps\sum_{k=1}^{n-1}\tilde \theta^{n-k}V_k^N,
\end{align*}
with $V_1^N=U_1^N=\|h_N^*-h_{\eps,N}^{*_1}\|_{W^{1,1}}$. In other words, 
\begin{align*}
  V_n^N&=KC\frac{\ln(N)}{N}\frac{\tilde \theta}{1-\tilde \theta}+K\eps\sum_{k=1}^{n-1}\tilde \theta^{n-k}V_k^N\\
  &=KC\frac{\ln(N)}{N}\frac{\tilde \theta}{1-\tilde \theta}+K\eps \tilde \theta V_{n-1}^N +\tilde \theta K \eps\sum_{k=1}^{n-2}\tilde \theta^{n-1-k}V_k^N\\
  &=KC\frac{\ln(N)}{N}\frac{\tilde \theta}{1-\tilde \theta}+K\eps \tilde \theta V_{n-1}^N +\tilde \theta \left( V_{n-1}^N-KC\frac{\ln(N)}{N}\frac{\tilde \theta}{1-\tilde \theta}\right)\\
&=KC\frac{\ln(N)}{N}\tilde \theta+(K\eps +1)\tilde \theta  V_{n-1}^N.
\end{align*}
Consequently, choosing $\eps_{3}^*$ such that $(K\eps_{3}^* +1)\tilde \theta:=\tilde\theta_1<1$, for any $\eps <\eps_{3}^*$ the sequence $(V_n^N)_{n\in \mathbb N}$ is of the form
\begin{align*}
V_n^N:=\tilde\theta_1^{n-1}V_1^N+KC\frac{\ln(N)}{N}\frac{\tilde \theta}{1-\tilde \theta}.
\end{align*}
Thus, using the above in \eqref{eq:fixpoint}, for any $n\in \mathbb N$,

\begin{align*}
\|h_{\varepsilon,N}^*-h_\eps^*\|_{W^{1,1}}
&\leq C_K \gamma^n+\tilde\theta_1^{n-1}\|h_N^*-h_{\eps,N}^{*_1}\|_{W^{1,1}}+KC\frac{\ln(N)}{N}\frac{\tilde \theta}{1-\tilde \theta}.\\
\end{align*}
One can then choose $n=-\ln(N)(\ln(\gamma)\vee \ln(\tilde\theta_1))$ to obtain
\begin{align*}
\|h_{\varepsilon,N}^*-h_\eps^*\|_{W^{1,1}}&\leq  C_K N^{-1}+N^{-1}\|h_N^*-h_{\eps,N}^{*_1}\|_{W^{1,1}}+KC\frac{\ln(N)}{N}\frac{\tilde \theta}{1-\tilde \theta}\\
&\leq  C_K N^{-1}+N^{-1}KC\frac{\ln(N)}{N}+KC\frac{\ln(N)}{N}\frac{\tilde \theta}{1-\tilde \theta}.
\end{align*}
\end{proof}

\end{document}
